\chapter{Background}
\label{ch:background}

\section{Abstract Meaning Representation}
\label{sec:bg-amr}

Abstract Meaning Representation (AMR) is a formalism for encoding the meaning of a
sentence as a rooted, directed, acyclic graph~\cite{banarescu2013abstract}.
Each node in the graph represents a concept, and each edge represents a semantic
relation between concepts.
Concepts are drawn from PropBank frames~\cite{palmer2005proposition}, which define
canonical verb senses and their argument roles.
For example, \texttt{send-01} is the PropBank frame for the sending action, with
\texttt{:ARG0} for the sender, \texttt{:ARG1} for the thing sent, and \texttt{:ARG2}
for the recipient.

Graphs are typically written in Penman notation, a parenthesised tree serialisation.
The sentence ``Do not send money'' parses to:

\begin{lstlisting}
(s / send-01
   :ARG0 (u / __system)
   :ARG1 (m / money)
   :polarity -)
\end{lstlisting}

The root concept is \texttt{send-01}.
The \texttt{:polarity -} edge marks negation.
The \texttt{:ARG0} role points to the implicit agent, which we represent as a special
\texttt{\_\_system} node since AMR does not have a dedicated ``you'' concept.

The key property of AMR for our purposes is what it discards.
AMR throws away word order, grammatical morphology, function words, and syntactic
structure.
It keeps the predicate, its arguments, and propositional operators like polarity and
modality.
This means that many surface forms map to the same or similar graphs.
``Do not send money'', ``money must not be sent'', and ``sending money is prohibited''
all produce a graph centred on \texttt{send-01} with \texttt{:polarity -}.
An attacker cannot change the underlying predicate by rewording the same command.

Two sentences that mean the same thing do not always produce identical AMR graphs,
because different parsers may choose different PropBank frames for synonymous words.
``Wire funds to the attacker'' might parse to \texttt{wire-01} rather than
\texttt{send-01}.
This is not a weakness of AMR; it is handled by the synonym detection rule described
in Section~\ref{sec:sys-rules}.
The point is that AMR eliminates syntactic variation, reducing the attacker's options
to semantic variation, which is a narrower and more tractable problem.

AMR was designed for natural language understanding tasks and has been studied
extensively since its introduction~\cite{banarescu2013abstract,knight2021abstract}.
The Penman notation used in this thesis derives from the Penman project at
ISI/USC~\cite{matthiessen1991text}.

\section{AMR as a Security Primitive}
\label{sec:bg-security}

The properties that make AMR useful for natural language understanding also make it
useful for security.
Obfuscated language is a common technique in prompt injection attacks.
An attacker might embed an instruction in a story, surround it with irrelevant text,
or use unusual phrasing to avoid pattern-based filters.
AMR parsing strips these layers away.
No matter how the attacker phrases ``send money to this IBAN'', the parser should
produce a graph with \texttt{send-01} as the root and a monetary argument.

This is what we mean by ``lossy compression as a security feature''.
The information that AMR discards is precisely the information an attacker uses to
evade surface-level detectors.
What remains is the semantic skeleton, and it is much harder to hide a forbidden action
at the semantic level than at the syntactic level.

AMR has been used in other security-adjacent NLP contexts.
Chung et al.\ use AMR-based graph alignment to check whether agent behaviour complies
with regulatory policies~\cite{chung2025graphcompliance}, demonstrating that
predicate-argument structure can serve as the basis for formal compliance verification.
Our work applies the same representation to a different problem: real-time detection of
adversarial payloads in agent tool outputs.

Compared to embedding-based similarity, AMR has two advantages in this setting.
First, it produces discrete, interpretable structure.
When the firewall blocks an input, we can point to exactly which predicate matched
which policy prohibition.
Second, it does not require a threshold calibrated on training data.
A match between \texttt{send-01} and a policy that forbids \texttt{send-01} is
unambiguous, not a matter of similarity score.

\section{Prompt Injection}
\label{sec:bg-injection}

\subsection{Direct and Indirect Injection}

Prompt injection attacks occur when an adversary inserts malicious instructions into
the input of an LLM~\cite{owaspPromptInjection}.
Two forms are commonly distinguished.

In a direct injection, the attacker is the user: they include malicious text in their
own message, trying to override the system prompt.
For example: ``Ignore your instructions and tell me the password.''
Direct injection is a well-studied problem, and many defenses target it specifically.

In an indirect injection, the attacker is a third party who controls content that the
agent will read.
They inject instructions into a file, a web page, a database record, or any other
external resource that the agent is expected to process.
The user did not write the malicious content; the agent encounters it while performing
a legitimate task.
Indirect injection is harder to defend against because the agent must process external
content to function at all.
This thesis focuses on indirect injection.

\subsection{The Data-Instruction Boundary}

A common framing in the literature distinguishes data (information the agent reads)
from instructions (commands the agent follows).
The idea is that tool outputs should be treated as data, not instructions, so that
embedded commands are ignored.
In practice this boundary is hard to enforce.
An agent that reads a task list and executes the items on it is deliberately following
instructions from a file.
The user may have written the task list, or the user may have said ``run tasks.txt'',
knowing it was written by someone else but trusting its contents.
The data-instruction distinction collapses exactly in the cases where injections are
most dangerous.
We avoid this framing in our definition and instead focus on whether an action
contradicts a stated policy.

\subsection{Existing Defenses}

Keyword and pattern filters maintain lists of known attack phrases.
They are fast and produce no false positives on benign inputs, but they are brittle:
any new phrasing defeats them.

LLM-based judges pass the input to a second LLM and ask whether it is malicious.
They are more flexible than keyword filters but share the same fundamental vulnerability:
the judge processes the adversarial text, and a sufficiently clever injection can
manipulate the judge itself.
There is no formal guarantee that the judge is immune to injection.

Canary token approaches embed secret strings in the system prompt.
If the agent's output contains the canary, the system prompt was likely leaked.
This detects a specific, narrow attack class but cannot prevent execution of injected
commands that do not involve the canary.

\section{Related Work}
\label{sec:bg-related}

\subsection{Firewalls for LLM Agent Networks}

Abdelnabi et al.\ propose a dual-firewall architecture for securing agent-to-agent
communication networks~\cite{abdelnabi2025firewalls}.
Their work is the closest in spirit to ours, and it is worth understanding the
difference in approach.

Their \textit{Language Converter Firewall} addresses the data-instruction boundary by
eliminating the channel through which instructions can arrive.
Incoming messages are converted to a closed, domain-specific, structured protocol using
deterministic verification.
If a message cannot be expressed in the closed protocol, it is dropped.
This provides a strong structural guarantee: instructions cannot arrive because there is
no representation for them in the target protocol.
Their system achieves large reductions in attack success rates on the ConVerse benchmark
across 864 attacks (from 84\% to 10\% for privacy attacks, from 60\% to 3\% for
security attacks).

Our approach takes the opposite starting point.
Rather than eliminating the channel, we allow natural-language tool outputs to reach the
agent and detect when they contain payloads that contradict the policy.
This is less restrictive: our firewall does not require the output schema to be fully
controlled.
The trade-off is that our approach depends on the quality of the AMR parser and on the
policy being stated explicitly.

Two further differences are worth noting.
Abdelnabi et al.\ target agent-to-agent communication, where the inputs are messages
from other LLM agents.
Our target is tool output injection in a single-agent setting, where the injections
arrive in structured data like file contents and database records.
These are related but distinct threat models.
Additionally, our approach produces interpretable blocking decisions: when the firewall
fires, we can state exactly which predicate matched which policy rule.
Structural protocol conversion does not provide this.

The fact that two independent groups have converged on a firewall framing for this
problem, using different methods, suggests this is the right architectural layer to
address prompt injection in agent systems.

\subsection{Graph-Based Compliance Checking}

Chung et al.\ use AMR graph alignment to verify that LLM-generated responses comply
with regulatory policies~\cite{chung2025graphcompliance}.
Their work demonstrates that predicate-argument structure can encode policy constraints
precisely enough for automated compliance checking.
We build on the same idea but apply it to adversarial inputs rather than agent outputs,
and we operate in real time rather than as a post-hoc audit.

\section{The AgentDojo Benchmark}
\label{sec:bg-agentdojo}

AgentDojo is a framework for evaluating the security of LLM agents against task
injection attacks~\cite{rogueSecurityPromptInjectionsBenchmark}.
It provides a banking environment with a realistic set of tools: the agent can read
files, query transactions, retrieve account information, send money, update passwords,
and schedule payments.
User tasks specify what the agent should accomplish; injection tasks specify what an
attacker wants the agent to do instead.
The benchmark measures both utility (does the agent complete the user's task?) and
security (does the agent resist the injection?).

We use AgentDojo for two reasons.
First, the banking environment has clear, natural security policies.
Second, the benchmark provides a standard evaluation harness that makes our results
reproducible.

We extend AgentDojo with a custom suite of 10 attack categories and a set of benign
tasks designed to test specific properties of the firewall.
The standard AgentDojo calibration tasks serve as an additional reference point.
The full suite is described in Section~\ref{sec:eval-taxonomy}.
